\section{Tikz}
\tikz \draw[->] (0,0) -- (2,1) -- (3,2);

\begin{tikzpicture}
    \draw (0,0) circle (1); %small circle
    \draw(2,0) circle (1.5in); %big circle
    \draw(5,0) ellipse (10pt and 20pt);

    \draw node at (3,0) {$f(x)$};

    \filldraw (6,2) circle (0.1cm) node[anchor=west]{Anchored Node};
\end{tikzpicture}

\vspace{1in}

\begin{tikzpicture}
    \draw (0,0) rectangle (5,4);
    \draw (0,0) grid (5,4);
\end{tikzpicture}

\vspace{1in}

\begin{tikzpicture}
    \draw[ultra thick] (0,3) -- (2,3);
    \draw[very thick] (0,2.5) -- (2,2.5);
    \draw[thick] (0,2) -- (2,2);
    \draw[thin] (0,1.5) -- (2,1.5);
    \draw[very thin] (0,1) -- (2,1);
    \draw[ultra thin] (0,0.5) -- (2,0.5);

    \draw node at (3,3) {Ultra Thick};
    \draw node at (3,2.5) {Very Thick};
    \draw node at (3,2) {Thick};
    \draw node at (3,1.5) {Thin};
    \draw node at (3,1) {Very Thin};
    \draw node at (3,0.5) {Ultra Thin};
\end{tikzpicture}

\newpage

\begin{center}
    \begin{tikzpicture}[transform canvas={scale=4.0}] %[scale=4] changes distances not the canvas
        \draw[blue] (0,1) arc (90:-90:0.5cm and 1cm);
        \draw[dashed, red] (0,1) arc (90:270:0.5cm and 1cm);
        \draw (0,0) circle (1cm);
        \filldraw[red] (0,1) circle (0.05); %add fill=, and draw= to have separate colours
        \filldraw[red] (0,-1) circle (0.05);
        \shade[ball color=blue!10!white, opacity=0.2] (0,0) circle (1cm); %\fill can be used to make the entire background the same colour. 
    \end{tikzpicture}
\end{center}

\vspace{2in}

\begin{tikzpicture}[SIR/.style={rectangle, draw=red!60, fill=red!5, very thick, minimum size = 5mm}] %Defines the style SIR
    %Nodes
    \node[SIR] (Susceptible) {Susceptible $S(t)$}; %creates Node with SIT style
    \node[SIR] (Infectious) [below=of Susceptible] {Infectious $I(t)$}; %Positions this node below node with Id Susceptible
    \node[SIR] (Recovered) [below=of Infectious]{Recovered $R(t)$};

    %Lines
    \draw[->, very thick] (Susceptible.south) to node[right] {$a$} (Infectious.north); %Draws an arrow from the bottom of the Susceptible node to the top of the Infectious node, with a label of 'a' to the right of the line
    \draw[->, very thick] (Infectious.south) to node[right] {$b$} (Recovered.north);
    \draw[->, very thick] (Recovered.east) .. controls +(right:7mm) and +(right:7mm)..(Susceptible.east);
\end{tikzpicture}

\vspace{2in}

\begin{tikzpicture}[
        youngnode/.style={rectangle, draw=red!60, fill=red!5, very thick, minimum size=40},
        oldnode/.style={rectangle, draw=blue!60, fill=blue!5, very thick, minimum size=40},
    ]
    %Nodes
    \node[oldnode]        (SusO)                            { $S_O(t)$};
    \node[oldnode]        (InfO)       [below=of SusO]      { $I_O(t)$};
    \node[oldnode]        (RecO)       [below=of InfO]      { $R_O(t)$};

    \node[youngnode]      (SusY)        [left=of SusO]      { $S_Y(t)$};
    \node[youngnode]      (InfY)        [left=of InfO]      { $I_Y(t)$};
    \node[youngnode]      (RecY)        [left=of RecO]      { $R_Y(t)$};

    %Lines
    \draw[->, very thick] (SusO.south east)  to node[right] {$a_{OO}$} (InfO.north east);
    \draw[->, very thick] (InfO.south)  to node[right] {$b_O$} (RecO.north);
    \draw[->, very thick] (RecO.east)  .. controls  +(right:17mm) and +(right:17mm)   .. (SusO.east);

    \draw[->, very thick] (SusY.south west)  to node[left] {$a_{YY}$} (InfY.north west);
    \draw[->, very thick] (InfY.south)  to node[left] {$b_Y$} (RecY.north);
    \draw[->, very thick] (RecY.west) .. controls  +(left:17mm) and +(left:17mm)   .. (SusY.west);

    \draw[dashed,->, very thick] (InfO.north west)  to  (SusY.south east);
    \draw[->, very thick] (SusY.south east)  to node[left] {$a_{OY}$} (InfY.north east);

    \draw[->, very thick] (SusO.south west)  to node[right] {$a_{YO}$} (InfO.north west);
    \draw[dashed,->, very thick] (InfY.north east)  to  (SusO.south west);
\end{tikzpicture}